% !TeX encoding = UTF-8
\documentclass[variant=guia,cover=institutional,mode=screen]{ua-engineering-v6-article}
% !TeX encoding = UTF-8
\UASetUniversity{Universidad Austral}
\UASetFaculty{Facultad de Ingeniería}
\UASetDepartment{Departamento de Computación}
\UASetCampus{Pilar}
\UASetCareer{Ingeniería Informática}
\UASetCommission{Comisión A}
\UASetProfessor{Equipo docente}
\UASetSemester{Segundo cuatrimestre 2026}
\UASetCourse{Sistemas Distribuidos}
\UASetDate{2026}

\UASetClassNumber{08}
\UASetClassTitle{Causalidad, consenso y consistencia distribuida}
\title{Clase 08 --- Guía de ejercicios}
\author{\UAProfessor}
\date{\UADate}
\begin{document}
\maketitle
\section{Propósito}
Esta guía transforma cada concepto de la clase en una ejecución que debe reconstruirse. No se evalúa la cantidad de nombres recordados, sino la coherencia entre problema, información disponible, falla, garantía, mecanismo, costo y evidencia.
\section{Ruta de uso}
\subsection{Ruta recomendada de uso}
La guía está pensada para tres pasadas. En la primera, resuelva sin consultar los apuntes y marque con un signo de interrogación cada supuesto que no pueda justificar. En la segunda, compare su timeline, conjuntos o genealogía con el material de clase y corrija únicamente aquello que pueda explicar. En la tercera, redacte una defensa breve que incluya garantía, falla, costo y evidencia. No convierta la guía en un glosario: cada respuesta debe reconstruir una ejecución.

\begin{itemize}
\item \textbf{Durante la clase:} 1, 2, 3, 4, 7, 8, 13, 16 y 20.
\item \textbf{Consolidación:} 5, 6, 9, 10, 11, 12, 14 y 15.
\item \textbf{Desafío/Trabajo Final:} 17, 18 y 19.
\end{itemize}
En ejercicios temporales incluya timeline; en quórums, conjuntos; en versiones, genealogía; en Spanner, intervalo y condición de espera.

\section{Formato esperado de cada respuesta}
Cada ejercicio debe resolverse con una estructura mínima común:
\begin{enumerate}
\item \textbf{Actores y estado inicial:} qué sabe cada nodo antes del primer mensaje.
\item \textbf{Operación e invariante:} qué intenta hacer el cliente y qué historia está prohibida.
\item \textbf{Falla o concurrencia:} timeout, crash, pausa, partición o escrituras simultáneas.
\item \textbf{Mecanismo:} regla de reloj, mayoría, quórum, versión, repair o espera.
\item \textbf{Frontera de respuesta:} evento exacto que vuelve legítimo el ACK o resultado visible.
\item \textbf{Costo y límite:} latencia, indisponibilidad, metadatos, conflicto o deuda operativa.
\item \textbf{Evidencia:} timeline, métrica, consulta, prueba de falla o contraejemplo.
\end{enumerate}
Una respuesta que solo nombra ``Raft'', ``CAP'' o ``Dynamo'' se considera incompleta. Debe explicar cómo cambia la observación del usuario al mover un mensaje o una falla. Cuando haya varias alternativas correctas, declare el supuesto que hace válida su elección y la condición bajo la cual la revisaría.

\section{Ejercicios}
\begin{uaexercise}
\textbf{1. Timeout e indistinguibilidad.} Buenos Aires envía una reserva a Madrid y recibe timeout. Enumere cuatro historias compatibles con esa observación. Para cada una, indique qué evidencia adicional permitiría distinguirla.
\end{uaexercise}

\begin{uaexercise}
\textbf{2. Happened-before.} Dibuje tres líneas de proceso ---BA, líder y Madrid--- para click, send, receive y append. Marque dos relaciones causales, una relación transitiva y un par concurrente.
\end{uaexercise}

\begin{uaexercise}
\textbf{3. Lamport.} Asigne timestamps de Lamport a la ejecución anterior. Justifique la regla $L\leftarrow\max(L,t_m)+1$ y construya un par con $L(a)<L(b)$ pero sin $a\rightarrow b$.
\end{uaexercise}

\begin{uaexercise}
\textbf{4. Relojes vectoriales.} Con componentes [BA, Madrid, Líder], calcule los vectores hasta que el líder recibe ambas reservas. Explique por qué el resultado del líder es [2,2,3] y no [2,2,5].
\end{uaexercise}

\begin{uaexercise}
\textbf{5. Máquina de estados.} Defina estados y transiciones para VIP-01. Incluya reserva de Sofía, reserva de Tomás y repetición idempotente de un comando ya aplicado.
\end{uaexercise}

\begin{uaexercise}
\textbf{6. Safety y liveness.} Formule dos propiedades de safety y dos condiciones de liveness para el cluster de cinco réplicas. Muestre una ejecución que conserva safety pero pierde liveness.
\end{uaexercise}

\begin{uaexercise}
\textbf{7. Raft y frontera de commit.} Complete los logs de N1--N5 cuando N1 propone la entrada 87 en term 12. Marque cuándo está local, replicada, committed, aplicada y confirmada.
\end{uaexercise}

\begin{uaexercise}
\textbf{8. Crash móvil.} Analice dos crashes del líder: uno después del append local y otro después de la mayoría pero antes del ACK. Indique qué puede responder el nuevo líder y qué debería observar Sofía.
\end{uaexercise}

\begin{uaexercise}
\textbf{9. Total-order broadcast.} Explique por qué una máquina de estados replicada necesita los mismos comandos en el mismo orden. Dé un contraejemplo donde timestamps Lamport no basten.
\end{uaexercise}

\begin{uaexercise}
\textbf{10. Pausa, GC y fencing.} N1 se pausa; N2 es elegido en un término mayor; N1 despierta y escribe sobre un recurso externo. Diseñe una protección mediante término y otra mediante fencing token.
\end{uaexercise}

\begin{uaexercise}
\textbf{11. Raft y Multi-Paxos.} Mapee term, index, leader y followers a ballot, slot, proposer y acceptors para la decisión 87. Identifique la abstracción común y una diferencia operacional.
\end{uaexercise}

\begin{uaexercise}
\textbf{12. Tres logs.} Para log de consenso, WAL local y log Kafka, indique productor, unidad, garantía, recovery y una confusión frecuente.
\end{uaexercise}

\begin{uaexercise}
\textbf{13. CAP sin slogan.} Durante una partición 3|2, Sofía ya fue committed en la mayoría y Tomás llega a la minoría. Compare responder éxito, rechazo y pendiente. Identifique qué propiedad se sacrifica.
\end{uaexercise}

\begin{uaexercise}
\textbf{14. PACELC y recencia.} Compare una lectura local de 15 ms posiblemente atrasada con una coordinada de 180 ms. Proponga dos operaciones que elijan políticas distintas y métricas para evaluarlas.
\end{uaexercise}

\begin{uaexercise}
\textbf{15. Consistencia.} Diferencie recencia, linearizabilidad, serializabilidad y consistencia externa mediante cuatro historias mínimas.
\end{uaexercise}

\begin{uaexercise}
\textbf{16. Mayoría vs quorum.} Compare mayoría Raft con N=5,R=3,W=3 en un sistema Dynamo-like. Explique por qué tres respuestas no significan lo mismo.
\end{uaexercise}

\begin{uaexercise}
\textbf{17. Límites de R+W>N.} Construya tres contraejemplos: escrituras concurrentes, sloppy quorum y last-write-wins con reloj desalineado. En cada uno explique qué intersección existe y qué garantía falta.
\end{uaexercise}

\begin{uaexercise}
\textbf{18. Versiones y merge.} Modele un carrito con dos versiones concurrentes y diseñe un merge. Luego aplique la misma estrategia a VIP-01 y explique por qué falla.
\end{uaexercise}

\begin{uaexercise}
\textbf{19. Dynamo y reparación.} Diseñe una ejecución con nodo preferido caído, sloppy quorum, hint, hinted handoff, read repair y anti-entropy. Defina cuatro métricas de deuda de reparación.
\end{uaexercise}

\begin{uaexercise}
\textbf{20. Spanner y TrueTime.} Una transacción multi-grupo recibe $TT.now()=[4,16]$ ms y el coordinador elige $s=16$ ms. Muestre cuándo puede hacerse visible, explique external consistency y analice qué cambia si aumenta la incertidumbre.
\end{uaexercise}
\section{Criterios de entrega}
Toda respuesta debe declarar supuestos, frontera de garantía, falla cubierta, costo desplazado y una observación que permitiría refutar la explicación.
\end{document}
